\section{Related Work}
\label{sec:related_work}

\medskip\noindent\textbf{Ensembles and answer aggregation.}
Classical ensemble methods show that simple aggregation rules can be strong baselines when component errors are at least partly decorrelated \citep{breiman1996bagging,freund1997decisiontheoretic,kittler1998combining,dietterich2000ensemble,zhou2012ensemble}.
That perspective matters directly here: a pooled four-answer ensemble over frontier, L1, S1, and TALE is a meaningful comparator rather than a straw baseline.
The paper therefore does not claim broad superiority over ensembling; instead, it asks whether interpretable trace-guided routing can be competitive with simple pooled aggregation under the same already-generated answer pool.

\medskip\noindent\textbf{Self-consistency and reasoning-time aggregation for LLMs.}
Chain-of-thought prompting and zero-shot chain-of-thought showed that explicitly eliciting reasoning steps can improve mathematical and symbolic problem solving \citep{wei2022chain,kojima2022large}.
Self-consistency then aggregates over multiple sampled reasoning paths by selecting the most frequent answer \citep{wang2023self}, with later work extending this idea to broader candidate pools \citep{chen2023universalselfconsistency}.
Tree-structured and graph-structured reasoning methods broaden inference-time exploration further by organizing candidate traces into explicit search procedures \citep{yao2023tree,besta2024graph}, while program-aided reasoning couples generation to executable intermediate structure \citep{gao2023pal}.
Failure-Trace Allocator is adjacent to this line of work but different in scope: it does not sample new reasoning paths at selection time and it does not choose by answer frequency alone.
Instead, it selects among already generated answers using runtime-visible trace metadata plus external agreement signals.

\medskip\noindent\textbf{Verifier-based reranking and iterative correction.}
Another line of work uses learned verifiers, reward models, or self-correction mechanisms to rerank or revise candidate solutions.
For GSM8K, \citet{cobbe2021training} showed that an outcome verifier can materially improve answer selection when trained with correctness supervision.
Process reward models refine this idea by scoring intermediate reasoning steps rather than only final answers \citep{lightman2024lets}.
Related work on self-verification and iterative self-correction asks the model to critique or refine its own reasoning before committing to a final answer \citep{weng2023largelanguagemodelsbetter,madaan2023selfrefineiterativerefinementselffeedback}.
The method studied here differs in two ways: it uses deterministic gold-free routing predicates rather than a learned verifier, and it does not invoke additional corrective generation after the candidate pool is available.

\medskip\noindent\textbf{Selective prediction, deferral, and calibration.}
Selective prediction studies when a system should abstain, defer, or route away from its default prediction on subsets where reliability is lower \citep{geifman2017selectiveclassificationdeepneural,geifman2019selectivenet}.
Calibration analyses are relevant because such routing decisions depend on whether confidence-like signals carry meaningful reliability information \citep{guo2017calibrationmodernneuralnetworks}.
Failure-Trace Allocator is closer to this selective-routing perspective than to pure majority voting: the frontier answer remains the default on most examples, while LDG and ECO intervene only on a narrow runtime-identifiable subset.
This framing also clarifies the paper's proposition-level claim: if a gold-free runtime signal isolates examples where an alternative selector is more reliable than the frontier default, then conditional deferral on that subset can improve overall accuracy.

\medskip\noindent\textbf{Budgeted inference and test-time compute.}
Recent work on test-time compute asks how additional inference effort should be allocated to improve reasoning accuracy \citep{snell2024scaling,wu2024inference}.
Budget-forcing approaches such as s1 regulate the amount of reasoning performed during generation \citep{muennighoff2025s1}, while broader cost-aware systems such as FrugalGPT and language-model cascades study routing under resource constraints \citep{chen2023frugalgpt,dohan2022language}.
Adaptive computation inside a model or computation graph provides a related perspective on conditional compute allocation \citep{graves2016adaptive}.
The present paper addresses a narrower problem: after candidate-producing methods have already run under the same per-method budget cap, can a deterministic zero-extra-call selector use logged trace metadata to improve the final answer choice?
